% --- Idioma y Codificación ---
\usepackage[utf8]{inputenc}
\usepackage[spanish]{babel}

% --- Matemáticas Avanzadas ---
\usepackage{amsmath, amssymb, amsthm} % Esenciales para física
\usepackage{bm} % Para negritas en símbolos matemáticos (vectores)
% --- Gráficos y Colores ---
\usepackage{graphicx}
\graphicspath{{img/}} % Indica donde están las imágenes
\usepackage{xcolor}

% Definición de tus colores personalizados
\definecolor{navyblue}{RGB}{0, 0, 128}      % Azul Marino
\definecolor{darkred}{RGB}{150, 0, 0}
% Entorno personalizado para Ejemplos en Azul
\newenvironment{ejemplo}[1][]{
	\itshape\color{navyblue}
	\medskip
	\noindent\textbf{Ejemplo #1:} 
}{
	\medskip
}

% Comando rápido para texto importante en Rojo
\newcommand{\importante}[1]{{\color{darkred}\textbf{#1}}}

% --- Código de Programación (Perfecto para tus scripts de Python/C++) ---
\usepackage{listings}
\definecolor{codegreen}{rgb}{0,0.6,0}
\lstset{
	language=Python,
	basicstyle=\ttfamily\small,
	keywordstyle=\color{blue},
	commentstyle=\color{codegreen},
	breaklines=true,
	frame=single
}

\usepackage{amsthm}
\newtheorem{definition}{Definición}[section] % Numeración por sección

% --- Hipervínculos ---
\usepackage[colorlinks=true, linkcolor=blue, urlcolor=cyan]{hyperref}
